\startchapter[title={Introducción}, reference=Capitulo_1]

%----------------------------------------------------------------------------------------
\startsection[title={Problema de pickeo de órdenes}, reference=sec:problema]

En un centro de distribución, el {\em picking} (la recolección de los ítems
que componen un pedido dentro de un almacén) es una de las actividades que
más tiempo y costo insume dentro de las operaciones logísticas, llegando a
representar hasta el 50-75\% del costo operativo total de un almacén
\cite[urzua2019evaluating]. Un caso real documentado en un centro de
distribución nacional muestra que el tiempo de preparación de un pedido podía
llegar a las 8 horas, y que con una mejor asignación de las ubicaciones de picking
es posible reducir significativamente ese tiempo \cite[gutierrez2019reduction].
Reducir el tiempo de desplazamiento de los operarios tiene un impacto
económico directo y escalable, ya que cualquier mejora se repite en cada
turno, en cada operario y en cada pedido.

Llamamos {\bf OPP} (del inglés {\em Order Picking Problem}) al problema
de determinar, dado un conjunto de pedidos de clientes agrupados en lotes, a
qué operario ({\em picker}) asignar cada lote, en qué orden secuenciarlos, y
la ruta que cada operario debe seguir dentro del almacén para recolectar
todas las ubicaciones requeridas por sus lotes, de forma tal de minimizar una
función objetivo que combina el tiempo total de recolección con penalidades
por adelanto y demora respecto de la fecha de entrega comprometida
({\em due date}) de cada pedido. El almacén se representa como un grafo
dirigido cuyos nodos son las ubicaciones de picking, el depósito de partida y
el punto de entrega, y cuyos arcos representan los desplazamientos posibles
entre ubicaciones junto con su distancia o tiempo asociado.

El agrupamiento de pedidos en lotes ({\em batching}) es un problema
relacionado, pero distinto: en este trabajo los lotes se reciben ya
conformados como dato de entrada, y no forman parte de lo que el modelo
decide.

\stopsection

%----------------------------------------------------------------------------------------
\startsection[title={Enfoque propuesto}, reference=sec:enfoque]

Existen dos grandes familias de enfoques para resolver problemas de
optimización combinatoria como el OPP: los métodos {\bf exactos}, que
garantizan encontrar la solución óptima pero cuyo tiempo de cómputo crece
rápidamente con el tamaño de la instancia, y los métodos {\bf heurísticos}
(y metaheurísticos), que renuncian a la garantía de optimalidad a cambio de
tiempos de resolución mucho menores y de mayor escalabilidad. En un contexto
operativo real, donde las decisiones de ruteo deben tomarse en segundos y los
almacenes pueden tener cientos de ubicaciones, esta disyuntiva entre calidad
de la solución y tiempo de cómputo es central.

Este trabajo recorre esa disyuntiva en cuatro pasos sucesivos, cada uno
construido sobre el anterior:

\startitemize[n]
	\item {\bf Modelo de referencia.} Se toma como punto de partida el
	modelo de \ac{mip} propuesto por Coruzzolo et al.
	\cite[coruzzolo2023orderpicking], que resuelve de forma conjunta la
	asignación, la secuenciación y el ruteo de los lotes.
	\item {\bf Modelo monolítico.} Con ese modelo como guía, se formula un
	modelo lineal propio que incorpora dos características del problema real
	que la formulación de referencia no contempla: la velocidad heterogénea
	entre operarios y los requisitos de cadena de frío. El mismo recibe el nombre de
	{\em monolítico} porque decide ruta, asignación y secuenciación en un
	único modelo.
	\item {\bf Modelo en dos etapas (Split).} Para reducir los tiempos de
	ejecución del modelo monolítico, se lo descompone en dos modelos que se
	resuelven uno después del otro: el primero calcula la ruta más corta de
	cada lote, y el segundo, tomando esos tiempos de ruta como datos, decide
	la asignación, la secuenciación y la cadena de frío. A este modelo lo
	llamamos {\em Split}.
	\item {\bf Metaheurísticas.} Como el modelo Split sigue siendo
	inviable en tiempos prácticos para instancias moderadamente grandes, se
	programan dos metaheurísticas, basadas en \ac{hc} y en \ac{aco},
	{\em equivalentes en características} a los modelos lineales: resuelven
	el mismo problema, con la misma función objetivo, las mismas
	restricciones de velocidad y cadena de frío, y sobre el mismo espacio de
	soluciones, pero sin garantía de optimalidad. Esto permite comparar sus
	resultados directamente contra el óptimo.
\stopitemize

Los modelos lineales se desarrollan en el Capítulo~\in[Capitulo_3] y las
metaheurísticas en el Capítulo~\in[Capitulo_4].

\stopsection

%----------------------------------------------------------------------------------------
\startsection[title={Objetivo}, reference=sec:objetivo]

\startsubsection[title={Objetivo general}]

Evaluar, sobre un conjunto de escenarios comparables, el compromiso entre
calidad de la solución, tiempo de cómputo y robustez que ofrecen distintos
enfoques de resolución (exacto y heurísticos) al problema de asignación,
secuenciación y ruteo de picking, para caracterizar en qué condiciones cada
enfoque resulta más apropiado.

\stopsubsection

\startsubsection[title={Objetivos específicos}]

\startitemize
	\item Formular e implementar, tomando como guía la formulación de
	referencia, un modelo lineal (monolítico) que incorpore velocidad
	heterogénea entre operarios y requisitos de cadena de frío.
	\item Optimizar el modelo monolítico descomponiéndolo en un modelo lineal
	en dos etapas (Split) que alcance el mismo valor óptimo en menor tiempo.
	\item Diseñar e implementar dos metaheurísticas, basadas en \ac{hc} y en
	\ac{aco}, equivalentes a los modelos lineales en las características que
	modelan, de forma tal que sus resultados sean comparables entre sí.
	\item Diseñar un conjunto de escenarios de prueba (funcionales y de estrés)
	que permitan comparar los cuatro modelos bajo condiciones equivalentes.
	\item Medir y comparar, sobre esos escenarios, la tasa de convergencia, el
	tiempo de ejecución y la calidad de la solución respecto del óptimo.
	\item Identificar limitaciones estructurales de los enfoques heurísticos
	evaluados y proponer líneas de trabajo futuro para superarlas.
\stopitemize

\stopsubsection
\stopsection

%----------------------------------------------------------------------------------------
\startsection[title={Organización del trabajo}, reference=sec:organizacion]

El trabajo se organizó en seis etapas.

En la etapa de {\bf investigación y análisis} se relevó la literatura
sobre OPP y sobre los métodos exactos y heurísticos aplicables, y se
estudió en detalle el modelo de referencia de Coruzzolo et al.
\cite[coruzzolo2023orderpicking] para identificar qué características del
problema real (velocidad heterogénea entre operarios, requisitos de cadena de
frío) no estaban contempladas en esa formulación.

En la etapa de {\bf modelado} se formuló y programó, con el modelo de
referencia como guía, el modelo monolítico, que lo extiende con esas
características.

En la etapa de {\bf optimización} se aplicaron distintas estrategias con la
finalidad de reducir los tiempos de procesamiento del modelo monolítico, que
culminaron en su descomposición en dos etapas, dando lugar al modelo lineal
{\em Split}.

En la etapa de {\bf desarrollo de metaheurísticas} se programaron las
metaheurísticas basadas en \ac{hc} y \ac{aco}, equivalentes en
características a los modelos lineales, y se refinaron con la finalidad de
pulir los resultados obtenidos y mejorar el tiempo de convergencia al óptimo.

En la etapa de {\bf experimentación} se diseñó un conjunto de escenarios de
prueba, funcionales y de estrés, y se ejecutaron los cuatro modelos sobre esos
escenarios midiendo tasa de convergencia, tiempo de ejecución y calidad de la
solución respecto del óptimo.

Finalmente, en la etapa de {\bf comparación} se analizaron esos resultados
para caracterizar el compromiso entre calidad de la solución, tiempo de
cómputo y robustez que ofrece cada enfoque, y para identificar limitaciones
estructurales de las heurísticas evaluadas.

\stopsection
\stopchapter
